%
% teil2.tex -- Beispiel-File für teil2 
%
% (c) 2020 Prof Dr Andreas Müller, Hochschule Rapperswil
%
% !TEX root = ../../buch.tex
% !TEX encoding = UTF-8
%
\section{Der Weierstraß-Produktsatz}
\kopfrechts{Der Weierstraß-Produktsatz}

\subsection{Das Problem: Unendliche Produkte und ihre Konvergenz}

Bei einem Polynom bestimmen die Nullstellen $a_1, \ldots, a_n$ das Polynom bis
auf einen konstanten Vorfaktor eindeutig:
\[
p(z) = c \cdot (z - a_1)(z - a_2) \cdots (z - a_n).
\]
Für holomorphe Funktionen mit unendlich vielen Nullstellen stellt sich die Frage,
ob sich eine solche Darstellung als unendliches Produkt verallgemeinern lässt.
Dabei tauchen zwei grundlegende Schwierigkeiten auf.

\subsubsection{Normierung der Faktoren}
Statt der Faktoren $(z - a_n)$ verwendet man zweckmäßigerweise Faktoren der Form
$\bigl(1 - z/a_n\bigr)$.
Der Unterschied ist mehr als kosmetisch: Beim Vorfaktor $c$ eines Polynoms
steckt die Information über den Betrag und die Phase bei einem festen Punkt;
bei einem unendlichen Produkt würde ein schlecht normierter Faktor zu einem
divergenten Gesamtvorfaktor führen.
Die Wahl $\bigl(1 - z/a_n\bigr)$ normiert jeden Faktor so, dass er bei $z = 0$
den Wert $1$ annimmt, was die Konvergenzanalyse entscheidend vereinfacht.

\subsubsection{Unbestimmtheit durch ganze Funktionen}
Bei Polynomen bestimmen die Nullstellen das Polynom \emph{bis auf einen
konstanten Faktor} $c \in \mathbb{C}^*$, also bis auf $e^b$ mit $b \in \mathbb{C}$.
Bei holomorphen Funktionen ist die Situation komplizierter: Die Nullstellen
bestimmen die Funktion nur bis auf einen Faktor $e^{g(z)}$, wobei $g$ selbst
eine beliebige ganze Funktion sein kann, nicht nur eine Konstante.
Der Exponent $g(z)$ ist gewissermaßen die ``transzendente Freiheit'', die über
die algebraische Information der Nullstellenmenge hinausgeht.

\subsubsection{Konvergenz des unendlichen Produkts}
Das Produkt $\prod_n (1 - z/a_n)$ konvergiert genau dann, wenn die Summe
$\sum_n z/a_n$ konvergiert. Diesen Zusammenhang sieht man am einfachsten durch
Logarithmieren: Der Logarithmus eines Produkts ist die Summe der Logarithmen,
und für kleine $x$ gilt $\log(1 - x) \approx -x$. Genauer:
\[
\log(1 - x)
=
-\biggl(x + \frac{x^2}{2} + \frac{x^3}{3} + \cdots\biggr).
\]
Mit $x = z/a_n$ ergibt sich also, dass die Konvergenz des Produkts im Wesentlichen
von der Konvergenz der Reihe $\sum_n 1/a_n$ abhängt.

Dies ist jedoch keineswegs garantiert. Wählt man zum Beispiel $a_n = n$,
so divergiert $\sum_n 1/n$ als harmonische Reihe, das naive Produkt
$\prod_n (1 - z/n)$ konvergiert also nicht.

\subsubsection{Weierstraß-Faktoren als Lösung}
Die Idee zur Behebung dieses Konvergenzproblems ist elegant: Man betrachtet
den störenden ersten Term $x$ in der Reihenentwicklung von $\log(1-x)$ und
addiert ihn wieder hinzu:
\[
\log(1 - x) + x
=
-\biggl(\frac{x^2}{2} + \frac{x^3}{3} + \frac{x^4}{4} + \cdots\biggr).
\]
Für kleines $x$ ist $x^2/2$ viel kleiner als $x$, sodass die entsprechende Summe
$\sum_n (z/a_n)^2/2$ deutlich bessere Chancen hat zu konvergieren.
Indem man die Exponentialfunktion nimmt, bedeutet dies: Man ersetzt den Faktor
$(1 - x)$ durch den modifizierten Faktor $(1-x)\,e^{x}$. Gerade das ist der
\emph{Weierstraß-Faktor} $E_1(x) = (1-x)\,e^{x}$  \cite{produkt:weierstrass}.

Reicht dies immer noch nicht aus wenn also auch die $x^2/2$-Terme zu langsam
abklingen, so substrahiert man noch einen weiteren Term aus dem Logarithmus:
\[
\log(1-x) + x + \frac{x^2}{2}
=
-\biggl(\frac{x^3}{3} + \frac{x^4}{4} + \cdots\biggr).
\]
Der führende Term ist jetzt $x^3/3$, was die Konvergenz nochmals verbessert.
Exponenzieren ergibt den Faktor $E_2(x) = (1-x)\,e^{x + x^2/2}$.

Allgemein: Muss man $p$ störende Terme aus dem Logarithmus eliminieren, so
erhält man den \emph{Weierstraß-Faktor} $E_p(x)$.
\index{Weierstrass-Faktor@Weierstraß-Faktor}%
\index{Elementarfaktor}%
\index{Ep(x)@$E_p(x)$}%
Die Ordnung $\varrho$ einer ganzen Funktion bestimmt dabei, wie schnell die
Nullstellen wachsen, und daraus ergibt sich die notwendige Wahl von $p$.

\subsection{Elementarfaktoren und Konvergenz}

Das zentrale technische Hilfsmittel sind also die Weierstraß-Elementarfaktoren.
Für $p = 0, 1, 2, \ldots$ definiert man:
\begin{equation}
E_0(z) = 1 - z,
\qquad
E_p(z)
=
(1-z) \cdot \exp\!\biggl(z + \frac{z^2}{2} + \cdots + \frac{z^p}{p}\biggr)
\quad \text{für } p \ge 1.
\end{equation}
Der Exponentialterm ist so konstruiert, dass er genau die ersten $p$ Terme
in der Laurent-Reihe von $\log(1-z)$ kompensiert und damit die Divergenz
\index{Laurent-Reihe}%
des Faktors $(1-z)$ für wachsende Argumente unterdrückt.
Genauer gilt die Abschätzung
\begin{equation}
|1 - E_p(z)| \le |z|^{p+1}
\quad \text{für } |z| \le 1.
\end{equation}
Diese Ungleichung ist der Schlüssel für den Konvergenzbeweis des Produktsatzes;
für einen vollständigen Beweis sei auf \cite{produkt:conway}
oder \cite{produkt:steinShakarchi} verwiesen.

\subsection{Formulierung des Satzes}

Weierstraß bewies den folgenden fundamentalen Satz \cite{produkt:steinShakarchi}:

\begin{satz}[Weierstraß, 1876]
\index{Faktorisierungssatz!von Weierstraß}%
Sei $f$ eine ganze Funktion mit einer Nullstelle der Ordnung $m \ge 0$ bei $z = 0$
und weiteren Nullstellen $a_1, a_2, \ldots$ (mit Vielfachheiten,
nach wachsendem Betrag geordnet).
Dann existiert eine ganze Funktion $g(z)$ und nicht-negative ganze Zahlen $p_n$ mit
\begin{equation}
f(z)
=
z^m \cdot e^{g(z)} \cdot
\prod_{n=1}^{\infty} E_{p_n}\!\biggl(\frac{z}{a_n}\biggr).
\end{equation}
Das Produkt konvergiert gleichmäßig auf Kompakta.
Für ganze Funktionen endlicher Ordnung $\varrho$ genügt die konstante Wahl
$p_n = \lfloor \varrho \rfloor$.
\end{satz}

\subsection{Rückkehr zum Sinus}

Die Sinusfunktion $\sin(\pi z)$ ist eine ganze Funktion der Ordnung $\varrho = 1$,
da $|\sin(\pi z)| \le e^{\pi|z|}$ und die Nullstellen $a_n = \pm n$ linear wachsen.
Man wählt daher $p_n = 1$, also $E_1(z) = (1-z)e^z$.
Wegen der Symmetrie $\sin(\pi(-z)) = -\sin(\pi z)$ heben sich die Exponentialterme
der Paare $+n$ und $-n$ gegenseitig auf.
Dies weil $e^{z/n} \cdot e^{-z/n} = 1$ und so erhält man:
\begin{equation}
\sin(\pi z)
=
\pi z \cdot e^{g(z)} \cdot \prod_{n=1}^{\infty} \biggl(1 - \frac{z^2}{n^2}\biggr).
\end{equation}

Eine Schärfung des Weierstraß-Satzes genannt Hadamard-Faktorisierungssatz existiert.
\index{Hadamard}%
\index{Faktorisierungssatz!von Hadamard}%
Dieser ist für Funktionen endlicher Ordnung definiert und impliziert, dass $g(z)$ ein Polynom vom Grad
höchstens $\lfloor \varrho \rfloor = 1$ sein muss, also $g(z) = Az + B$.
Aus $\sin(\pi z)/(\pi z) \to 1$ für $z \to 0$ folgt $B = 0$,
und die Symmetrie, dass $\sin(\pi z)$ ungerade ist, ergibt $A = 0$.
Damit gilt Eulers Formel rigoros:
\begin{equation}
\sin(\pi z)
=
\pi z \cdot \prod_{n=1}^{\infty} \biggl(1 - \frac{z^2}{n^2}\biggr).
\end{equation}
